% ============================================================
%  第三章
% ============================================================
\section{GNO 的观察与 Langlands 对偶}
\markboth{第三章\ GNO 的观察与 Langlands 对偶}{第三章\ GNO 的观察与 Langlands 对偶}

\noindent{\small\color{shuyiGray}对应底本：原 p.13--p.18。}\par\medskip

\srcp{13}接着，GNO 做出了那个引起 Michael Atiyah 兴趣的观察：对任何单 Lie 群 $G$，都存在另一个单 Lie 群 $\gvee$，其极大环面记为 $\Tvee$，使得
\[
  \Hom(T,U(1))=\Hom(U(1),\Tvee),
\]
并且反过来，
\[
  \Hom(U(1),T)=\Hom(\Tvee,U(1)).
\]
换句话说，在 GNO 的解释里（至少当两边的规范对称性都被约化、或被“希格斯化”到一个极大环面时），$G$ 的电电荷就是 $\gvee$ 的磁荷，反之亦然。当然，你们中有些人会注意到：正如 Atiyah 最初指出的那样，GNO 的对偶群 $\gvee$ 与 Langlands 对偶群是同一个东西。

\srcp{14}Montonen 与 Olive 第二篇论文的要点，是要为这一切在量子动力学里安排一个位置。

\subsection{量子规范理论与耦合常数}

\srcp{15}量子规范理论是经典规范理论的推广，它依赖于一个额外的参数，即规范耦合 $e$（对单规范群而言只有一个参数）。当 $e$ 很小时，量子规范理论退化为经典规范理论；当 $e$ 不小时，量子规范理论相当难以理解。

\srcp{16}规范群不同的规范理论在经典上并不等价。但 Montonen 与 Olive 提出，它们在量子力学意义上可以等价。更具体地说，他们提出了一个具体的例子：一对规范理论，一个的规范群是 $G$，另一个的规范群是 $\gvee$；他们主张这两者在量子力学意义上等价，而在经典意义上不等价。这种等价应当把第一个理论里的电电荷映到第二个理论里的磁荷，反之亦然。因为这个缘故，Montonen--Olive 对偶常被称为\emph{电—磁对偶}。对偶应当把耦合常数反演：
\begin{equation}
  e^{\vee 2}=\frac{4\pi}{e^2},
  \label{eq:coupling}
\end{equation}
其中 $e^\vee$ 是 $\gvee$ 理论的规范耦合。由于 $e$ 与 $e^\vee$ 不可能同时很小，这就解释了为什么对偶在经典层面看不出来。

\subsection{\emph{Terra incognita}}

\srcp{17}要恰当地传达 Montonen--Olive 猜想在当时显得多么狂野，是很难的。量子规范理论当时只在相当小的 $e$ 下被理解，而大 $e$、量子效应很大的区域，真正是一片未知之地（\emph{terra incognita}）。Montonen--Olive 猜想看起来像是一跃，不仅跳出了我们已经理解的范围，而且跳出了我们可能理解的范围。但无论如何，Montonen 与 Olive 有一些引人注目的可能证据：他们证明，某些关于各种粒子（带电荷和/或磁荷的粒子）质量的经典公式，满足一些关系；这些关系若在量子力学意义上成立，就会从他们的对偶猜想推出。

\begin{yizhu}[terra incognita]
\textit{terra incognita} 是拉丁语“未知之地”，旧地图上用来标注未探明的区域。
原文用它形容大耦合区域，语气比“未知领域”重；译文保留原文并加引号，不改成
“未开垦的处女地”一类会引入别的联想的说法。
\end{yizhu}


\subsection{伦敦：两个疑虑}

\srcp{18}我听从 Atiyah 的建议，安排去伦敦找 Olive 谈，但等我到那里的时候，我已经相当怀疑了。让我不安的是 GNO 那篇论文的一点：电与磁的匹配只有在希格斯化之后——也就是约化到极大环面之后——才成立。我觉得一个根本性的陈述不应当依赖于这种希格斯化。让我不安的是 Montonen--Olive 那篇论文的另一点：他们在技术上假定“希格斯势”为零，而按一切标准逻辑，这个假定在量子力学上是不自洽的。

\begin{pizhu}{根数据：对偶的全部内容}
设 $T\subset G$ 是极大环面，$X^*(T)=\Hom(T,U(1))$ 是\emph{特征标格}，$X_*(T)=\Hom(U(1),T)$ 是\emph{余特征标格}。二者有自然的整值配对
\[
  \ip{\chi}{\lambda}=\chi\circ\lambda\in\Hom(U(1),U(1))=\ZZ .
\]
$G$ 的\emph{根数据}是四元组 $(X^*(T),\Phi,X_*(T),\Phi^\vee)$，其中 $\Phi\subset X^*(T)$ 是根，$\Phi^\vee\subset X_*(T)$ 是余根。\emph{对偶根数据}把两个格互换，同时把根与余根互换：
\[
  (X^*(T),\Phi,X_*(T),\Phi^\vee)\ \longmapsto\ (X_*(T),\Phi^\vee,X^*(T),\Phi).
\]
由这个对偶根数据构造出的单连通群就是 $\gvee$。于是
\[
  X^*(\Tvee)=X_*(T),\qquad X_*(\Tvee)=X^*(T),
\]
把两边代入 $\Hom$ 就得到第 13 页的两个等式。用 Cartan 矩阵说：$A_{ij}=\ip{\alpha_i}{\alpha_j^\vee}$，对偶就是矩阵转置；这正是根系统的对偶把 $B_n$ 与 $C_n$ 互换、而 $F_4$、$G_2$ 保持不变的原因。

\smallskip
\noindent\textbf{常见的对偶对：}
\begin{center}\small
\begin{tabular}{@{}ll@{}}
\toprule
$G$ & $\gvee$ \\
\midrule
$SU(N)$ & $PSU(N)=SU(N)/\ZZ_N$ \\
$Sp(n)$ & $SO(2n+1)$ \\
$Spin(2n)$ & $SO(2n)/\ZZ_2$ \\
$E_8,\ F_4,\ G_2$ & 自对偶 \\
\bottomrule
\end{tabular}
\end{center}
注意 $SU(N)$ 与 $PSU(N)$ 的李代数相同，对偶群却不同：对偶看到的不只是李代数，而是整个根数据（李代数加上整体结构）。这也是 GNO 论文的一个副产品——它给出了一个纯物理的判据，用来区分李代数相同而整体结构不同的群。
\end{pizhu}

\begin{pizhu}{“对偶群”的三种含义}
本章出现了三个容易混淆的“对偶”，值得当场分清。

\begin{enumerate}
  \item \textbf{GNO 对偶群。}纯粹由根数据定义，交换 $X^*$ 与 $X_*$。它是一个复（或紧）李群 $\gvee$，不依赖任何动力学。
  \item \textbf{Langlands 对偶群。}在数论的情形，还要把 Galois 群（或 Weil 群）的作用加进来，记作 ${}^LG=\gvee\rtimes W_F$。对分裂的情形、以及本书讨论的复曲线情形，它就退化为 $\gvee$。所以 Witten 说“同一个东西”，是在分裂情形下说的，这一点在第 7 章会再出现。
  \item \textbf{电磁对偶中的“对偶理论”。}指以 $\gvee$ 为规范群、耦合为 $e^\vee$ 的那个量子场论。它与原来的理论是同一个理论的两种描述；这里的“对偶”是一个关于动力学的断言，不是群论。
\end{enumerate}

三者共用符号 $\gvee$ 并非偶然：第 3 条之所以可能，正是因为第 1 条。但读者若把三者当成同一件事，就会在“为什么 $SU(2)$ 的对偶是 $SO(3)$”和“为什么 $SU(2)$ 理论的对偶还是 $SU(2)$ 理论”之间卡住。
\end{pizhu}

\begin{pizhu}{耦合常数、$\tau$ 与那个 $4\pi$}
式 \eqref{eq:coupling} 的系数依赖 $e$ 的归一化，值得交代清楚。今天通行的写法是把 $\theta$ 角与耦合合成
\[
  \tau=\frac{\theta}{2\pi}+\frac{4\pi\ii}{g^2},
\]
于是 $S$ 对偶就是 $\tau\mapsto-1/\tau$；取 $\theta=0$，得 $g'^2=(4\pi)^2/g^2$。原文写的是 $e^{\vee2}=4\pi/e^2$，与这个通行写法差一个 $4\pi$ 的因子。这不是笔误，而是 1977 年文献对 $e$ 的归一化与今天不同（当时常把 $4\pi$ 吸收进耦合或场的定义里）。

不变的说法只有一句：\emph{对偶把耦合取倒数}。凡是要引用具体系数的地方，都应当先声明归一化；本书正文照录原文，批注里的公式一律采用上面的 $\tau$ 约定。
\end{pizhu}
